\chapter{Never sign blind}

Never sign blindly. Prefer the hardware wallet screen
and independently checked calldata over any web UI.
A compromised front end can look perfect and still feed
you a drain.

The key is not stolen. You handed the spend over.

\section{Ground truth}

The laptop lies. The phone lies. The dApp lies. The
simulation can lie. The trusted display on the hardware
wallet is what you are authorizing. If that screen and
the website disagree, reject.

Demand clear signing. If the wallet wants a blind hash,
you are accepting a lot of risk. Sometimes a contract
forces it. Then the question is whether this interaction
is worth that.

\section{Simulate, then doubt it}

Run a simulator on anything that is not a boring
transfer you already understand. Read the predicted
balance changes. Then remember a simulation can show one
thing and execute another. It is a filter, not a
blessing.

\section{Four checks on an EOA}

\begin{enumerate}
  \item \textbf{Origin.} The URL. Clones are cheap.
  \item \textbf{Contract.} The ``interacting with''
  address, from official docs or a labeled explorer
  page, not from chat.
  \item \textbf{Function and parameters.} The name, the
  \texttt{to} / \texttt{recipient} / \texttt{onBehalfOf}
  field, the decimals.
  \item \textbf{Fee.} A simple transfer with a huge gas
  bill is a smell.
\end{enumerate}

\section{Practice when it is free}

Wise Signer and similar drills exist so the first time
you read calldata is not with rent money on the screen.
Use them.

The detailed Safe and Squads verification flows are in
the Protocol Multisig companion and on the site. The
rule does not change there: hash, simulate, decode, then
sign. Last signer executes only after that.
